\documentclass[11pt]{article}
\usepackage[a4paper,margin=1in]{geometry}
\usepackage{amsmath,amssymb,amsthm,microtype}
\newtheorem{theorem}{Theorem}[section]
\newtheorem{corollary}[theorem]{Corollary}
\theoremstyle{remark}\newtheorem{remark}[theorem]{Remark}
\title{Uniform Derivative Envelopes for Projection-Free $\det_2$ Tails}
\author{Bin Wang}\date{July 2026}
\begin{document}\maketitle
\begin{abstract}
The modulus-complete complement of RH-282 has an all-order geometric trace
envelope.  We show that the same estimate controls every fixed derivative of
the omitted logarithmic determinant tail.  If
$|\tau_{\sigma,n}|\le M_\sigma q^{n-2}$,
$M_\sigma\le C\sigma^{-\alpha}$, and $qR<1$, then the $s$th derivative of
$\sum_{n\ge m}\tau_{\sigma,n}z^n/n$ on $|z|\le R$ is bounded by a constant
times $M_\sigma m^{s-1}(qR)^m$.  A logarithmic clock strictly above the sharp
RH-283 threshold therefore gives convergence in every fixed $C^s$ norm.  In
the Hardy-disk instance the gain is a positive power of noise.  This theorem
controls only the complementary spectral factor; the finite head-to-
monodromy bridge remains open.
\end{abstract}

\section{Derivative estimate}
Assume
\[
 |\tau_{\sigma,n}|\le M_\sigma q^{n-2},\qquad n\ge2,
\]
and define
\[
 E_{\sigma,m}(z)=\sum_{n\ge m}\frac{\tau_{\sigma,n}}n z^n.
\]

\begin{theorem}\label{thm:derivative}
Let $x=qR<1$ and fix an integer $s\ge0$.  There is a finite constant
\[
 B_{s,x}=\sum_{k\ge0}(1+k)^{\max(s-1,0)}x^k
\]
such that, for $m\ge\max(2,s)$,
\begin{equation}\label{eq:bound}
 \sup_{|z|\le R}|E_{\sigma,m}^{(s)}(z)|
 \le M_\sigma q^{-2}R^{-s}B_{s,x}
 m^{\max(s-1,0)}x^m.
\end{equation}
\end{theorem}
\begin{proof}
For $s=0$, use $1/n\le1/m$.  For $s\ge1$, differentiating $z^n/n$ gives a
coefficient no larger than $n^{s-1}R^{n-s}$.  Write $n=m+k$ and use
$(m+k)^{s-1}\le m^{s-1}(1+k)^{s-1}$.  The remaining series is exactly
$B_{s,x}$.
\end{proof}

\begin{corollary}[smooth tail convergence]
If $M_\sigma\le C\sigma^{-\alpha}$ and
$m_\sigma=\lceil a\log(1/\sigma)\rceil$ with
$a\log(1/(qR))>\alpha$, then
\[
 E_{\sigma,m_\sigma}\longrightarrow0
 \quad\text{in }C^s(|z|\le R)
\]
for every fixed $s$.  The canonical-product tail
$\exp(-E_{\sigma,m_\sigma})$ converges to $1$ in the same topology.
\end{corollary}
\begin{proof}
The right side of \eqref{eq:bound} is a fixed power of
$\log(1/\sigma)$ times
$\sigma^{a\log(1/(qR))-\alpha}$.  It tends to zero.  The exponential claim
follows from the chain rule on a bounded neighborhood of zero.
\end{proof}

\section{Certified instance}
For $\alpha=1$, $q=1/2$, $R=7/5$, and $a=4$, the exponent is
\[
 4\log(10/7)-1=0.426699775\ldots>0.
\]
Thus the high-order noisy spectral factor from RH-282 tends to one together
with every fixed derivative on the full target disk.

\begin{remark}
Derivative control of the tail cannot identify the finite head.  In
particular it does not show that the modulus-selected noisy roots lie on the
finite monodromy shell.  Gates A--E remain false/open, and no arithmetic or
Hilbert--Polya conclusion is drawn.
\end{remark}
\end{document}
